\chapter{Moral Judgment: Psychological Mechanisms, Learning and Normative Assessment}
\label{sec:p11-rationality-12-3}

Wiegmann and Sauer's 2021 account asks two related but different questions:
how people generate moral judgments, and whether those judgments and the
processes generating them deserve rational confidence. Its descriptive
part compares dual-process, universal-moral-grammar and model-free/model-
based accounts. Its normative part follows the movement from developmental
rationalism through intuitionism and sentimentalism to the possibility of
rational moral learning. The authors' central proposal is that sensitivity
to morally relevant information, including the capacity to learn and revise,
is more informative than an unqualified emotion-versus-reason opposition
\parencite[pp.~701--707]{WiegmannSauer2021MoralJudgment}.

This does not make successful learning equivalent to moral truth. A process
can accurately learn a local convention, a reward contingency or a group's
prejudice. A study can predict acceptance of a sacrificial act without
measuring impartial concern or demonstrating the act's permissibility.
The source's note~1 explicitly warns that studies discussed for their
normative relevance were not necessarily designed to answer normative
questions. Their measured task, psychological interpretation and proposed
moral conclusion must remain separate.

\begin{hbmethod}{def:p11-r123-targets}{Four targets for evaluating a moral-judgment finding}
Specify the participant's response and the rule coding it, the psychological
mechanism proposed to generate it, the statistical or instrumental standard
used to assess that mechanism, and the moral standard used to assess the
judgment's content. Record the premises connecting these targets.

An acceptance response is not itself evidence that a whole ethical theory
was adopted. A neural contrast is not a uniquely identified cognitive
process. Bayesian updating on supplied information is not proof that the
information encodes morally sound rules. Observed endorsement, epistemic
warrant and moral authority require different arguments. Empirical records
follow \cref{def:p19-evidence-record}.
\end{hbmethod}

\section{Trolley tasks and the dual-process proposal}
\label{sec:p11-r123-dual}

The source describes the familiar System~1/System~2 contrast: relatively
automatic, quick and often unconscious processing versus relatively
effortful, analytic and consciously controlled processing. These bundles
are proposed process characteristics, not a definition assigning every
fast judgment to emotion and every slow judgment to rationality. The
review also associates the first with evolutionarily older and the second
with younger processing; those characterizations are not independent
evolutionary evidence established by the dilemma studies.
The
particular Greene account maps automatic affective responses primarily
onto deontological judgments and controlled cost--benefit assessment
primarily onto consequentialist judgments. Other dual-process views make
different claims, including the social-intuitionist proposal discussed below
\parencite[p.~702]{WiegmannSauer2021MoralJudgment}.

\begin{hbexample}{def:p11-r123-trolleys}{Switch, Push and TrapDoor}
In Switch, diverting a runaway train saves five people but kills one on
the alternative track. In Push, stopping the train requires pushing a
sufficiently heavy person from a bridge into its path, saving the five but killing that
person. TrapDoor retains the person's use as a physical stopper but
replaces pushing with operating a switch that opens a trap door.

These are stipulated judgment scenarios, not instructions for real
emergencies. Their equal stipulated casualty totals do not make every
morally relevant feature equal. Switch does not require the victim's
death or physical use to save the five: if the person escapes the track,
the diversion still works. Push and TrapDoor use the person as the means
of stopping the train. Physical contact, the action performed and causal
role are therefore distinct contrasts
\parencite[pp.~702--704]{WiegmannSauer2021MoralJudgment}.
\end{hbexample}

Experimental practice often calls approving the intervention
``consequentialist'' and rejecting it ``deontological.'' The source's
note~2 warns that these are loose operational labels: many deontological
theories permit intervention in Switch. Genuine consequentialist commitment
also requires an appropriate outcome criterion and impartial concern,
not simply reduced aversion to harming someone. A participant can accept
one sacrifice for a selfish or unrelated motive. Consequently, an item
response cannot by itself classify the person's overall moral outlook
\parencite[p.~702 and note~2]{WiegmannSauer2021MoralJudgment}.

The reviewed support for Greene's account includes neuroimaging contrasts,
cognitive-load effects and response-time differences interpreted as
controlled override of an affective intuition. Its proposed normative
extension treats some deontological intuitions as crude alarms sensitive
to irrelevant features, while favoring flexible consequence assessment.
The authors retain four separate objections:
\begin{enumerate}
\item Neuroimaging and latency findings may be more ambiguous than the
      initial interpretation suggests; activation-to-process inference
      requires discrimination among alternative processes.
\item Some philosophical critics argue that the relevant moral argument
      does not depend on neuroscience or an evolutionary explanation at all.
\item An emotional or intuitive origin need not be an epistemic defect.
\item Even a reliably identified defect may concern the task response
      rather than deontological theory. Intuitiveness can reverse across
      moral contents, and sacrificial acceptance can lack impartial concern.
\end{enumerate}
The source also retains the dispute between critiques of sacrificial
``utilitarian'' measures and later process-dissociation defenses; it does
not present one unopposed verdict on those measures. A process explanation
needs an additional moral-relevance premise before it can debunk a moral
theory \parencite[p.~702]{WiegmannSauer2021MoralJudgment}.

\section{An original intuitiveness/content study}
\label{sec:p11-r123-kahane}

\begin{hbdefinition}{def:p11-r123-kahane}{Kahane and colleagues' behavioral and imaging task}
The original article was published online in March 2011 and in the 2012
issue. Sixteen healthy right-handed volunteers participated, nine female,
mean age 29.25, range 21--41. They were screened for prior neurological
or psychiatric illness; the main paper does not give a total number
screened out or a representative recruitment frame. This is a volunteer
sample, not a population survey of moral rationality.

Eighteen independent pilot judges supplied immediate unreflective responses.
Eight dilemmas with at least twelve pilot judges choosing the designated
deontological option were classed DI; ten with at least twelve choosing
the designated utilitarian option were classed UI. Most DI cases concerned
serious harm; many UI cases concerned other duties, such as not lying.
The labels are pilot-majority classifications, not individual scanner
participants' measured first intuitions.

In four approximately ten-minute sessions, dilemma order was randomized.
Each trial used three text screens, followed by a yes/no judgment on a
suggested action; button mappings were balanced across participants.
Difficulty was rated from 0 to 100. Response time ran from the question
screen to the button press, although relevant information had appeared
earlier. Each decision was classified by both dilemma type and the
participant's chosen content. The resulting four cells distinguish
intuitive/counterintuitive from designated utilitarian/deontological
responses \parencite[pp.~394--395]{KahaneEtAl2012IntuitivenessR123}.
\end{hbdefinition}

Participants chose the pilot-intuitive response more often in each type:
reported tests were \(t(15)=5.81,p<.001\) for UI and
\(t(15)=-4.16,p=.001\) for DI. Across dilemma types, the correlation
of numbers of counterintuitive responses was \(r=.14,p=.606\), and
of designated utilitarian responses \(r=-.14,p=.599\).
These small-sample nonsignificant correlations do not establish independence
or stable absence of a trait; they caution against assuming one item-family
measure transparently represents a general moral orientation
\parencite[p.~396]{KahaneEtAl2012IntuitivenessR123}.

The reported difficulty contrast favored greater difficulty for
counterintuitive responses, \(F(1,11)=24.95,p<.001\), whereas the
utilitarian/deontological contrast was \(F(1,11)=.06,p=.811\).
Neither the intuitiveness response-time contrast,
\(F(1,11)=.272,p=.612\), nor the content response-time contrast,
\(F(1,11)=.314,p=.586\), was significant. Thus perceived difficulty
and response time did not give identical evidence. The main paper's
comparison degrees of freedom vary between 15, 13 and 11; it does not
fully state the inclusion flow reconciling those denominators. Recruitment
of sixteen must not silently become sixteen complete observations for
every contrast. Figure~2 gives standard-error bars, not confidence
intervals for every displayed comparison
\parencite[pp.~396--397]{KahaneEtAl2012IntuitivenessR123}.

The imaging analysis found patterns associated with intuitiveness and
some patterns associated with content. Its common global threshold was
\(p<.001\), uncorrected, with a five-voxel minimum cluster; supplementary
tables separately marked clusters surviving familywise correction.
Counterintuitive responses involved rostral anterior cingulate activity,
but the authors consider several interpretations, including emotional
conflict and social valuation. The main study does not show that a unique
controlled moral-reasoning process was identified. Its discussion explicitly
notes the lack of longer counterintuitive response times and of some
expected higher-level deliberative-region activations
\parencite[pp.~395--400]{KahaneEtAl2012IntuitivenessR123}.

The crossed comparisons challenge a universal equation of content with
process. They do not prove a universal reversed equation. Response cells
depend on subjects' choices, and the dilemma families differ in duty,
severity and content; randomized order does not remove those item
differences. Small sample, pilot-based classification, incomplete analysis-
flow reporting and absent separate supplement verification constrain
generalization. No causal change in moral truth, overall competence or
philosophical commitment was measured. The source paper itself warns
that its ``utilitarian'' choices do not entail a utilitarian outlook
\parencite[pp.~395,399--401]{KahaneEtAl2012IntuitivenessR123}.

\section{Universal moral grammar and acquisition}
\label{sec:p11-r123-grammar}

Mikhail's universal-moral-grammar proposal differs from a purely
domain-general dual-process account. An innate moral faculty would
represent actions through moral categories and assign deontic status
from their structure. The chapter's supporting families are early
development of distinctions, apparent cross-cultural commonalities,
people's difficulty articulating operative rules, and means/side-effect
patterns that a simple contact-based contrast misses. None is a complete
demonstration of innate moral grammar by itself
\parencite[pp.~702--703]{WiegmannSauer2021MoralJudgment}.

The stronger argument is poverty of the stimulus: the moral rules are
allegedly too subtle to be recovered from the available learning input,
so morality-specific innate constraints are proposed. That inference needs
both the actual input and the learner's possible hypotheses and learning
mechanisms. Failure to explain a rule verbally does not establish that
it could not have been learned. Conversely, a learned simulation does
not establish that all relevant innate structure is absent. The rational-
learning alternative below supplies a particular mechanism to test,
not a declaration that the nativist question is closed.

\section{Action value, outcome value and model-based structure}
\label{sec:p11-r123-models}

\begin{hbdefinition}{def:p11-r123-models}{Model-free and model-based evaluation in this source}
Model-based evaluation uses a representation of environmental causal
structure to assess expected consequences of available acts. Model-free
evaluation uses learned values associated with acts or situations, without
explicitly reconstructing those consequences at each choice. Reinforcement
learning supplies ways of updating such values through experience and
prediction error. These are distinctions of representation, learning and
action selection, not synonyms for emotion and thought.

Under the proposed decision criterion, the model-based chooser favors
the act with greatest expected value, while the model-free chooser favors
the best learned value. Neither criterion by itself justifies the values
as moral values.

Stored values can be computationally economical and successful in a
stable environment. A changed contingency can make them obsolete until
relevant experience or other updating occurs. Model-based flexibility
depends on the model's accuracy and the resources to use it; having a
model does not guarantee an optimal or morally correct decision
\parencite[p.~703]{WiegmannSauer2021MoralJudgment}.
\end{hbdefinition}

The chapter presents Crockett and Cushman as independently developing
related frameworks that preserve a dual-system distinction while making
it computationally more specific. It reviews action/outcome-aversion
dissociations as initial support, not one homogeneous experiment with
common sample and effect. Its suggestion that this account was the main
or most promising contender belongs to the 2021 assessment, not a
timeless empirical ranking of every model
\parencite[pp.~701,703--704]{WiegmannSauer2021MoralJudgment}.

For Push versus Switch, learned aversion to pushing or harming someone
can differ from learned value of operating a switch despite the same
stipulated saved/killed totals. That explanation needs an action encoding
and relevant learning history; the word ``model-free'' alone does not
supply either. For Switch versus TrapDoor, the physical action is now
similar. The source instead distinguishes whether harming the person
is a required subgoal in the causal plan. A represented aversive subgoal
can be evaluated like an action, producing an additional negative value
in the means case.

The two related accounts are not identical here. Note~3 records Crockett's
use of a third, Pavlovian system interacting with model-based structure,
whereas the preceding exposition describes model-free values assigned
to represented subgoals. This architectural difference must survive a
comparison. A domain-general system does not dispense with the need
to encode the morally relevant causal structure
\parencite[pp.~703--704, note~3]{WiegmannSauer2021MoralJudgment}.

The chapter reports an approval ordering of Switch above TrapDoor above
Push, without providing one common sample, effect magnitude or interval.
It also uses action/omission asymmetries: under its selected account,
aversive learned action values may not attach in the same way to the
always-available option of not acting. That is a model explanation of
a response asymmetry, not a proof that omissions lack moral consequences
or cannot themselves be actions under another description
\parencite[pp.~703--704]{WiegmannSauer2021MoralJudgment}.

\section{From developmental rationalism to intuitionist challenges}
\label{sec:p11-r123-history}

The normative discussion starts with the developmental tradition associated
with Piaget, Kohlberg and Turiel. The reviewed Kohlberg sequence progresses
from preconventional external sanctions through conventional rules and
roles to postconventional universal principles. This is a theory relating
development and sophistication, not a score assigning every child a
globally valid moral level. The source retains early moral/conventional
discrimination findings and Gilligan's objection that privileging abstract
rights and justice can marginalize morally significant relationships and
care. These challenges concern the task and evaluative criterion, not
merely the age at which a response appears
\parencite[p.~704]{WiegmannSauer2021MoralJudgment}.

Haidt's social-intuitionist model instead emphasizes rapid intuitive
judgment and often post-hoc reasons used for persuasion. Moral
dumbfounding describes maintaining a moral judgment while failing to
give an adequate articulated justification after proposed reasons have
been challenged. A failure to verbalize reasons is not automatically
absence of every reason or evidence that no revision is possible.
The source also reviews moral-foundations and political-difference work,
with challenges to how disagreements and their resolution are construed.
Its cited cases include liberal/conservative comparisons and
self-identified libertarian samples; those labels do not identify a
uniform population, task or moral outlook.
Those heterogeneous reports do not supply a representative political
population estimate or a universal mechanism in this chapter
\parencite[pp.~704--705]{WiegmannSauer2021MoralJudgment}.

Sentimentalist accounts concern not only the proximal production of a
particular judgment but the distal emotional foundations of valuing.
The reviewed argument invokes affective impairments and emotion
manipulations to claim necessity and sufficiency. The source immediately
retains challenges to both interpretation and robustness: incidental
affect effects may be small, and competence in moral/conventional
discrimination can survive some affective impairment. A clinical label,
one discrimination task, moral behavior and the capacity for extended
agency are not interchangeable outcomes. The chapter does not provide
their primary diagnostic, recruitment and uncertainty records and does
not authorize a general clinical verdict about a group
\parencite[pp.~704--705]{WiegmannSauer2021MoralJudgment}.

The rationalist replies are likewise several proposals. One treats
responsiveness to reasons as part of what makes a judgment moral. Another
objects to replacing temporally extended agency with isolated responses
to unusual dilemmas. Sauer's educated-intuitions proposal lets earlier
rational processes become automatized; Kumar's hybrid proposal treats
belief and motivation/emotion as clustering components. These are not
one experimentally established process. Each requires evidence suited
to the claimed timescale and construct
\parencite[p.~705]{WiegmannSauer2021MoralJudgment}.

Experimental-philosophy examples supply a further direction of influence:
evaluative judgments can affect causal attribution, intentional-action
attribution and judgments of personal identity. The source's Knobe and
related examples concern moral cognition affecting ostensibly nonmoral
cognition, not only emotion preceding moral judgment. Whether such
influence is legitimate depends on the attribution task and its standard;
it is not settled by naming every influence bias or by describing humans
as moralizers \parencite[p.~705]{WiegmannSauer2021MoralJudgment}.

\section{Moral intuitions and the rational-learning proposal}
\label{sec:p11-r123-learning}

The source follows Haidt in describing moral intuition as an assessment
suddenly appearing with affective valence, without conscious awareness
of intervening search, weighing or inference. This concerns access to
the generating steps, not proof that no informative computation occurred.
The possibility of override remains. Conversely, an apparently articulate
explanation does not establish that it generated the original response
\parencite[p.~705]{WiegmannSauer2021MoralJudgment}.

Computational work on reinforcement and Bayesian learning supplies an
alternative to treating intuition as inherently irrational. Sophisticated
learning can track subtle regularities without conscious access. The
source brings this into dialogue with moral epistemology rather than
assuming that every philosophical view requires one robust mind-independent
truth conception. It cites different defenses of retaining moral beliefs;
those are not collapsed into one psychological result
\parencite[pp.~705--706]{WiegmannSauer2021MoralJudgment}.

Railton's reinterpretation of the taboo sibling-incest vignette questions
whether the apparent removal of adverse outcomes also removes ex-ante
risks. A stipulated fortunate ending need not establish that taking the
risk was permissible. This is an alternative reading of what participants
may respond to in a vignette stipulating consent, contraception, secrecy
and a benign ending. These stipulations do not themselves provide an
empirical estimate of every ex-ante risk. The chapter supplies neither
an independently measured risk nor a settled moral verdict on the
scenario here. The source retains objections
to the dumbfounding interpretation from Jacobson, Guglielmo and Royzman
and colleagues. It also invokes affective information about others'
intentions and likely reactions, alongside lesion evidence, as possible
resources for moral decision making. Neither affective benefit nor an
impairment case proves that all emotion is reliable or all judgment
without it is impossible \parencite[p.~706]{WiegmannSauer2021MoralJudgment}.

\begin{hbdefinition}{def:p11-r123-dice}{The source's dice illustration under an explicit sampling law}
Choose uniformly one fair die with face set \(\{1,\ldots,d\}\), where
\(d\in D=\{4,6,8,10\}\), and retain it for \(n\) independent
rolls, with \(n\) a nonnegative integer. Let \(H_d\) denote the selected-die hypothesis and \(x\) the
ordered observed sequence. Put \(L_d(x)=\Pr(x\mid H_d)\) and
\(p_d(x)=\Pr(H_d\mid x)\), with the latter defined only when
the sequence has positive marginal probability.
The chapter's ten outcomes are
\[
                    x=(2,2,1,4,3,1,2,2,3,4).
\]
Equal prior selection, fairness and independent repeated rolls are the
declared premises. They are not assumptions automatically justified for
social instruction \parencite[p.~706]{WiegmannSauer2021MoralJudgment}.
\end{hbdefinition}

\begin{hbtheorem}{thm:p11-r123-size-principle}{proved=annex:p11.r123.size-principle}{The finite dice posterior and its scope}
If all observations lie in \(\{1,\ldots,d\}\), then
\(L_d(x)=d^{-n}\); otherwise it is zero. For the ten stated rolls,
every candidate is compatible and
\[
       p_d(x)=\frac{d^{-10}}{4^{-10}+6^{-10}+8^{-10}+10^{-10}}.
\]
For \(d=4,6,8,10\), respectively, the posterior values are approximately
.981910, .017028, .000959 and .000103. Among two compatible candidates
\(d<e\), their equal-prior posterior odds are \((e/d)^n\), favoring
the smaller candidate for \(n>0\). Empty data retain the prior. An
observation outside a candidate's faces eliminates it instead of favoring
it merely because it is small, provided some candidate can generate the
whole sequence. If none can, this posterior is undefined.
\end{hbtheorem}

The formula supplies the chapter's promised exact-probability illustration.
It is a conditional likelihood result, not a universal instruction to
prefer narrow rules. Unequal priors, selective presentation, dependent
samples or a different hypothesis space can change the inference.
Learning a word such as dog from selected examples and learning a conduct
rule both require those sampling and representation assumptions to be
specified. The chapter's child-learning analogy does not itself measure
them. Its actual word examples distinguish dogs from the wider animal
category, and Dalmatians from dogs generally: a few examples within a
narrow compatible category can favor that narrower extension under the
size-principle assumptions. This is not a guarantee under arbitrary
teaching or sampling.
\parencite[p.~706]{WiegmannSauer2021MoralJudgment}.

\section{An original rule-learning record and its limiting results}
\label{sec:p11-r123-nichols}

Nichols and colleagues' nested rule scopes distinguish intended outcomes,
foreseen side effects, agent-caused outcomes and outcomes in general.
A prohibition applying only to what an agent does is narrower than one
also requiring prevention of an outcome caused by others or nature.
Their own note warns that doing/allowing does not simply coincide with
act/omission: one can produce an outcome by omission. Intending also
need not imply confidently foreseeing success. The nesting is a model
of relevant interpretations with acknowledged boundary cases
\parencite[pp.~535--537]{NicholsEtAl2016RulesR123}.

The paper's CHILDES analysis used child-directed speech for four children,
Abe, Adam, Ross and Sarah, at 33--36 months. Independent coders identified
conduct-relevant instruction and assessed compatibility with narrow or
wide scope; reported agreement exceeded 99 percent, with disagreements
settled by discussion. More than 99 percent of coded communications were
consistent with narrow scope. That is compatibility of a corpus's
instruction, not observation that children adopted the narrow rule or
that the rule is morally justified. Four children's available records
do not establish all cultures' learning input
\parencite[pp.~537--538]{NicholsEtAl2016RulesR123}.

\begin{hbdefinition}{def:p11-r123-nichols-one}{Study 1: adult foreign-rule induction}
Twenty-four Mechanical Turk adults completed a long survey for nominal
cash payment; eighteen, eleven female, entered the reported analysis.
The paper excludes fast completions using wording about a threshold
below the 75th percentile of completion rate. That wording does not
transparently reconcile the ordering/threshold with retaining eighteen
of twenty-four. The reported count and criterion are preserved, not
silently repaired. Its footnote says the significance pattern is unchanged
when all participants are included; no independent reanalysis is supplied.

Participants inferred arbitrarily named rules in unemotional chalkboard,
shelf and litter domains. Trials supplied one intended-action example,
three intended-action examples, or three mixed examples including two
nonaction cases. Participants then classified additional intended and
overall-consequence cases as violations. These repeated binary rule-
extension judgments are not ratings of sacrificial moral permissibility.
Age distribution, recruitment country, collection dates and full ordering/
counterbalancing details are not given in the selected methods
\parencite[pp.~538--539]{NicholsEtAl2016RulesR123}.
\end{hbdefinition}

The paper reports five cross-domain errors among 1296 response opportunities,
not 1296 independently recruited people. Three narrow examples favored
narrow interpretation, reported \(t(17)=7.9,p<.0001\); mixed examples
favored a wider extension, \(t(17)=20.4,p<.0001\), with a reported
between-trial contrast \(t(17)=19.96,p<.0001\). Crucially, one versus
three narrow samples did not differ significantly,
\(t(17)=1.31,p=.20\). The authors expressly say this first study
does not establish the size principle over a subset alternative that
prefers the smallest compatible scope without a sample-number effect
\parencite[pp.~539,542--543]{NicholsEtAl2016RulesR123}.

Post-task similarity judgments helped construct a hierarchical hypothesis
space, and a Bayesian simulation used that space to resemble the response
pattern. A learned representation and qualitative fit are relevant
mechanistic evidence, but do not uniquely identify a Bayesian cognitive
process, constitute an out-of-sample model comparison or prove a child's
developmental trajectory \parencite[pp.~540--542]{NicholsEtAl2016RulesR123}.

\begin{hbdefinition}{def:p11-r123-nichols-two}{Study 2a: sample number under a changed task}
A separate Mechanical Turk sample comprised 56 completers, with eight
excluded for failing either of two comprehension checks; 48 remained,
15 female. The restaurant task described a teacher choosing one or three
violation clips from a listed set of eleven, involving placing napkins
on a windowsill or seeing and leaving them there. The test asked how
likely a leaving case and a placing case violated the learned rule,
using a 1--7 rating from not at all to extremely likely.

This changed the domain, strengthened instructions and replaced binary
classification with a graded likelihood judgment. The method does not
state participant counts by condition or a complete random-assignment
procedure. Teacher selection also differs from the uniform random
sampling stipulated in the dice example
\parencite[pp.~543--544]{NicholsEtAl2016RulesR123}.
\end{hbdefinition}

For the leaving case, reported means were 3.56 after one sample and 2.26
after three, \(t(46)=2.2,p<.05\). The 1.30 difference is on the
seven-point rating scale, not a percentage-point difference in moral
accuracy. For the placing case, means were 5 and 5.95,
\(t(46)=1.49,p=.142\). The latter nonsignificance is not an
equivalence result. The selected report does not give confidence intervals
for these mean contrasts or a pre-registration record. Expectations about
the teacher's informative selection, the instructions and the changed
response scale remain possible contributors
\parencite[p.~544]{NicholsEtAl2016RulesR123}.

The two studies constrain an account of adult rule induction and provide
a learning alternative to a blanket poverty-of-stimulus claim. Together
with the corpus they support a conditional account: appropriately structured
learners receiving suitably patterned input can acquire narrow-scope rules.
They do not demonstrate that all children do so, that every intended/
foreseen or doing/allowing distinction is learned by this mechanism, or
that narrow rules are morally superior. The chapter's claim that it
could be statistically irrational to infer broad utilitarian rules is
therefore relative to the proposed input, hypothesis space and sampling
law, not a general ethical refutation of utilitarianism
\parencite[p.~706]{WiegmannSauer2021MoralJudgment}.

\section{Learning reliability is not a moral safeguard}
\label{sec:p11-r123-normative}

The authors make room for optimism: intuitive judgments can be products
of sophisticated learning rather than crude impulses, and can potentially
respond to relevant information. They also retain Gjesdal's and Greene's
skeptical response. Domain-general learning can acquire irrelevant or
morally objectionable information; a welfare-sensitive learner can be
well adapted to personal advantage yet fail an impartial moral standard.
Statistical success at extracting an environment's regularities does not
guarantee that the regularities are ethically adequate
\parencite[pp.~706--707]{WiegmannSauer2021MoralJudgment}.

An improved social environment might make learning more morally useful,
for example through cooperation or contact that changes prejudices.
The review offers those possibilities, not a measured intervention effect
or a guarantee about past or future environments. Greene distinguishes
familiar within-group coordination from between-group conflicts, novel
institutional/technological problems and philosophical improvement. His
limits concern input, feedback and transfer, not a claim that no intuition
can ever be informative. Note~4 points to further moral-learning work,
while note~5 links the possible welfare/morality conflict to the distinct
Fehige--Wessels treatment
\parencite[pp.~706--707, notes~4--5]{WiegmannSauer2021MoralJudgment}.

The source's final conclusion is conditional: investigate when internal
learning mechanisms can obtain morally relevant external information.
There is no built-in safeguard selecting only morally relevant input or
only morally adequate output. Neither ``intuitive'' nor ``learned'' is
therefore an automatic verdict of irrationality, rational warrant or
moral authority. Feedback can be informative about what is approved
without being informative about what is right.

\section{Interfaces and bounded consultation}
\label{sec:p11-r123-use}

The distinction in \cref{sec:p19-rationality-11-2} between deontic
reasoning and deontological ethics prevents another shortcut: using
ought/permission vocabulary does not classify a response as a particular
moral orientation or cognitive process. A child's successful rule-violation
task is not the same measure as trolley acceptance, foreign-rule induction
or the legitimacy of a new norm. Their task benchmarks and interpretations
must be stated separately.

The morality/rationality comparison at \cref{sec:p11-rationality-12-1}
asks whether independently plausible practical and moral criteria converge.
Smith's \cref{sec:p11-rationality-12-2} asks which account of reasons and
rationally assessable intrinsic desires a psychology presupposes. The
present chapter investigates mechanisms and their epistemic relevance;
none of its empirical patterns proves those other chapters' normative
bridges. Broome's \cref{sec:p11-reasons-and-rationality} also keeps actual
reasons, motivating attitudes and successful response distinct. Learned
endorsement cannot replace those distinctions.

The developmental/rationalist claim that reflection contributes to moral
judgment is also different from Smith's Rationalism about a-priori moral
knowledge and moral requirements entailing reasons. The shared label
does not make a process result a resolution of that philosophical position.

Oxford's 2012 Chapter~19 by Waldmann, Nagel and Wiegmann is a distinct
source presentation cited by this chapter for the dilemma literature.
This treatment does not substitute the 2021 learning-centered account
for that complete 2012 scope. The shared subject and authorship do not
make their versions or coverage obligations identical.

\begin{hbmethod}{def:p11-r123-use}{Evaluate a mechanism and a moral claim without exchanging them}
Preserve the dilemma wording, response format, stage of judgment and
coding rule. State whether intuition means pilot-majority preference,
an individual's immediate answer or lack of access to intermediate
processing. Distinguish content, physical action, causal role, motivation
and the outcome standard.

For a process claim, identify what observations discriminate its rivals:
latency, difficulty, learning-input changes, selective outcomes or neural
contrasts are not interchangeable. Retain samples, exclusions, repeated
responses, analysis denominators, uncertainty and negative results. A
reported nonsignificant contrast is not proof that two processes are equal.

For a learning claim, state the hypotheses, priors, sampling mechanism,
feedback and task of transfer. Check whether a narrow-scope preference
could arise under a rival subset rule or a pedagogical interpretation.
For a moral conclusion, identify the standard and why the information
tracked is relevant to it rather than merely common, rewarded or
self-interested.

A successful application preserves Study 1's null sample-number contrast
while distinguishing Study 2a's positive result and task changes. Another
uses the dice posterior under its explicit probability law without treating
it as moral vindication. A failed application calls all sacrifice endorsers
utilitarian, interprets any brain activation as conscious ethical reasoning,
or infers a morally sound norm from successful acquisition of a local rule.
\end{hbmethod}

\begin{proofboundary}{def:p11-r123-boundary}{What the treatment establishes}
The source-specific account, examples, five notes and contested normative
bridges are reconstructed. The two primary empirical records retain their
actual methods and limits; wider cited studies remain attributed review
reports, without fabricated populations or measured effects. The finite
dice calculation has a full reciprocal annex proof. It proves no moral
axioms, universal psychological mechanism or ethical theory.

The original exclusion wording, analysis-denominator incompleteness and
unverified separate neuroimaging supplements remain explicit report
limitations. The chapter's model preferences and future-research hopes
remain dated to 2021. Empirical prediction, an explanation of response
generation and justification of moral content are not combined into
a universal moral-rationality score.
\end{proofboundary}

\index{moral judgment!psychological mechanisms}
\index{moral intuition!learning}
\index{trolley dilemmas!operational classifications}
\index{universal moral grammar}
\index{model-free learning!moral judgment}
\index{model-based learning!moral judgment}
\index{size principle!rule learning}
\index{moral dumbfounding}
